\section{Exclusion and completion of the proof}\label{sec:conclusion}

\subsection{The obstruction to a normal proof}

Here is the precise property of the ambient context used by the encoding.
Let $\mathcal F$ be its finite family of terminals.  Each $F\in\mathcal F$
is a beta-normal, sorted data expression containing only
$(\mathsf d,\mathsf d)$ labels.  A \emph{constant positive telescope} is
an expression
\[
 \Pi^{\mathsf d,\mathsf p}x_1:D_1.\,
 \cdots\Pi^{\mathsf d,\mathsf p}x_r:D_r.\,F,
 \qquad F\in\mathcal F,
\]
where $r\geq0$, the domains are normal data expressions, and each displayed
binder is absent from the subsequent suffix.  All the telescopes and suffixes
used below have specified sort typings.  Let $\mathcal T$ consist of the
converter domains, their suffixes, and all the terminals.

The ambient context is $\Delta=\Delta_{\mathsf d},\Delta_{\mathsf p}$.
Every variable in $\Delta_{\mathsf d}$ has data mode.  Every declaration
in $\Delta_{\mathsf p}$ is a proof variable of the form
\[
 c:\Pi^{\mathsf p,\mathsf p}u:T_c.F_c,
 \qquad T_c\in\mathcal T,\quad F_c\in\mathcal F.
\]
The domains $T_c$ are the constant positive telescopes just described.
All these types are normal.  Their typings, and the suffix typings, persist
in legal extensions of the context.

\begin{lemma}[No normal erasure at a terminal]
\label{lem:enc-no-normal-proof}
In the ambient context above, or in any legal extension by data declarations,
there is no proof-mode term $P$ with beta-normal erasure and a judgment
$P:T$ for any $T\in\mathcal T$.
\end{lemma}

\begin{proof}
Choose a counterexample of least syntax size, allowing all such data
extensions and all targets in $\mathcal T$.  A proof-mode term is neither
a sort nor a product.

Suppose first that $P$ is a lambda.  Its body and result have proof mode.
Its generated type is therefore a product with result label $\mathsf p$.
This type cannot convert to a terminal: a normal terminal is either not a
product, or has outer product labels $(\mathsf d,\mathsf d)$.
Consequently $T$ has a nonempty positive prefix.  Product compatibility
forces the lambda's labels to be $(\mathsf d,\mathsf p)$, with annotation
convertible to the first telescope domain.  Its binder is a data variable.
Generation gives a body typing whose expected type converts to the
remaining suffix.  The suffix's recorded sort typing, transported by
context conversion to the lambda annotation and weakened as necessary,
permits conversion of the body to that suffix.  The body has normal
erasure and is strictly smaller than $P$, in an allowed extension by one
more data declaration.  This contradicts minimality.

Otherwise $P$ is a variable or an application spine.  Every function position
of a proof-mode spine has proof mode.  Its head cannot be a sort or product,
and a lambda head with a nonempty spine would create a redex in its erasure.
Thus its head is a proof variable, hence one of the converters $c$.

If the spine is empty, variable generation says that $T$ converts to
$\Pi^{\mathsf p,\mathsf p}u:T_c.F_c$.  Both are normal.  The latter
has outer labels $(\mathsf p,\mathsf p)$, whereas $T$ either has a
$(\mathsf d,\mathsf p)$ prefix or is an all-data terminal.  This is
impossible by confluence and product compatibility.

If the spine is nonempty, let $Q$ be its first argument.  Generation at the
first application and at $c$, followed by product compatibility, forces
that application to have labels $(\mathsf p,\mathsf p)$ and its domain
to convert to $T_c$.  Thus $Q$ is a proof-mode term.  The recorded sort
typing of $T_c$, weakened to the current context, converts its generated
argument typing to $Q:T_c$.  Its erasure is normal because it is a subterm
of the normal erasure of $P$, and its syntax size is strictly smaller.
This is another counterexample to minimality.
\end{proof}

The lemma concerns proof mode, rather than every inhabitant of a terminal.
For example, the ambient data declaration $e_b:F_b$ is harmless: a data
variable cannot be the head of a proof-mode spine.  The reduction-lifting
lemma ensures that a constructed proof cannot change into such a data term
on the way to an erased normal form.

\begin{proposition}[Exclusion by system-wide weak normalization]
\label{prop:enc-exclusion}
Assume that every legal expression of the given pure type system is weakly
beta-normalizing.  For every primary component containing an odd closed walk and an
all-component profile triple, condition~\eqref{obs:forbidden} holds.
\end{proposition}

\begin{proof}
If the configuration prohibited in~\eqref{obs:forbidden} occurred,
Sections~\ref{sec:enc-wrappers}--\ref{rel-diagonal}, culminating in
Proposition~\ref{rel-contradiction}, would give a
finite legal labelled context $\Delta$ of the form above and a proof-mode
term
\[
                 \Delta\vdash P:F_b.
\]
Every other temporary assumption in the construction has been discharged
or instantiated.  Erasure gives an ordinary typing derivation
$|\Delta|\vdash |P|:|F_b|$ using only the original specification.
The system-wide hypothesis therefore applies to this particular open
legal term $|P|$.  Choose a finite reduction to a full beta-normal form.
Lemma~\ref{lem:enc-labelled-metatheory} lifts that reduction to a labelled
reduct $P'$ with the same exact type $F_b$, the same proof mode, and normal
erasure.  Lemma~\ref{lem:enc-no-normal-proof} rules out $P'$.
\end{proof}


\begin{proof}[Proof of Theorem~\ref{thm:main}]
By Lemma~\ref{pts:finite-restriction}, it suffices to treat a finite
specification. Process the finitely many primary components in dependency
order. For a component without an odd closed walk use
Proposition~\ref{prop:signed-component}. For one with no profile triple
entirely inside it use Proposition~\ref{prop:one-child-component}.
Every remaining component satisfies the hypotheses of
Proposition~\ref{prop:enc-exclusion}; hence
Proposition~\ref{obs:interface-theorem} supplies its candidate interface.
Proposition~\ref{prop:advance-component} therefore advances the
processed invariant at every stage. Once all components are processed,
Proposition~\ref{prop:interfaces-imply-sn} gives strong normalization
of every legal expression. This includes expressions in arbitrary valid
open contexts and every reduction position in their annotations.
\end{proof}
